\scanpage{68}
\mspage{\handwritten{- 3 -}}
2.3.1 et 4.3.1 pour $F$ parfait résultent de R-R quand on en dispose. Dire aussi un mot du degré d'un diviseur sur une courbe ?

p.26, li.2 : \deleted{\unsure{vider} ``quelconque''} remplacer ``annelé quelconque'' par ``commutativement annelé''.

p.27, li.9 : ``1.3.3'' et non ``1.2.3''.

p.27, li.5 du bas : dire ``passe au quotient par $\operatorname{Filt}^{2}(X)$ et induit ...''

p.29, li.4 : écrire plutôt $H^{\unsure{0|o}}(X,\underline{Z})$ \inserted{ou $\Gamma(\underline{Z}_X)$} au lieu de $\Gamma(\underline{Z})$.

\figure{\begin{tikzpicture}[x=0.03cm,y=-0.03cm,baseline]
\draw[line width=0.35pt] (2,13) .. controls (-2,23) and (4,24) .. (15,27) .. controls (24,30) and (25,40) .. (25,55);
\draw[line width=0.35pt] (25,58) -- (25,61);
\draw[line width=0.35pt] (190,0) .. controls (189,14) and (181,29) .. (167,34) .. controls (145,37) and (112,34) .. (86,35) .. controls (64,34) and (48,31) .. (37,34) .. controls (30,36) and (27,46) .. (26,56);
\end{tikzpicture}}

p.30, li.2 : il est peut-être prudent de supposer $X$ localement annelé.

p.30 : bas de page illisible.

p.31, li.1 : dire :``étant donné $M=T(\underline{O}^{\underline{n}_1},\ldots,\underline{O}^{\underline{n}_r})$ et la représentation ...''

p.31, li.4 : pourquoi supposer $M$ libre ? dire simplement ``par recollement de $M$ sur des cartes ...'' ou ``de la façon suivante.'', ce qui est \deleted{peut-être} encore mieux.

p.31, li.7 : ``tel que l'on ait ...''.Référer peut-être à la thèse de Giraud.

p.31, li.12 : supprimer ``fait opérer ... et'' et ``ainsi''. Après (6.1.4), dire qu'on \deleted{retrouve} a $M=T(\underline{O}^{\underline{n}_1},\ldots,\underline{O}^{\underline{n}_r})$.(référence Giraud). Pour (6.1.5), supposer $M$ localement libre.

p.32, \deleted{dire qu'on a une injection canonique} énoncer 6.2.2 en disant que l'injection canonique $\underline{\operatorname{Der}}(\operatorname{Gl}(\underline{n}))\handwritten{\unsure{\lhook\joinrel\relbar|\subset|\longleftarrow}}\operatorname{Sl}(\underline{n})$ est un isomorphisme. \handwritten{N'y a-t-il pas une référence dans Bible ou EGA3 ?}

p.34, li.10 du bas : ``voisinage étale'' au lieu de ``voisinage fpqc''.

p.\handwritten{3\unsure{\deleted{4}5|\deleted{5}4}}, li.12 du bas : même remarque.

p.35, li.4 du bas : que signifie $\operatorname{Gl}(n)(U)$ ? (si cela désigne les points de $\operatorname{Gl}(n)$ à valeurs dans $U$, alors la phrase n'a pas de sens).
% archive sheet number: 67
